\subsubsection{算子的例子}\label{subsec:operator-examples}

\Needspace{4\baselineskip}
\begin{Example}
\label{ex:dual-functional-bounded-operator}
对任意 Banach 空间 $X$, 由定义, 对偶空间 $X^*$ 的任意元素都属于 $\mathcal{L}(X,\mathbb{R})$.
\end{Example}

\Needspace{4\baselineskip}
\begin{Example}
\label{ex:identity-operator-and-inclusion}
当然, 由 $Ix=x$ 定义的从 $X$ 到 $X$ 的\textbf{恒等算子} $I$ 是有界算子的一个显然例子. 更有意义的是, 若 $X\subset Y$, 可以将从 $X$ 到 $Y$ 的恒等映射定义为把每个 $x\in X$ 映到同一 $x$ 的映射 (不过这次将其看作 $Y$ 的元素).
\end{Example}

以下引理是 Definition~\ref{def:compact-imbedding} 的直接推论.

\begin{Lemma}
\label{lem:identity-bounded-continuous-embedding}
设 $X$ 和 $Y$ 为 Banach 空间. 那么 $X$ 连续嵌入 $Y$ 的充分必要条件是从 $X$ 到 $Y$ 的恒等映射有良好定义 (即 $X\subset Y$) 且有界.
\end{Lemma}

在以下例子中使用这一点.

\Needspace{4\baselineskip}
\begin{Example}
\label{ex:sobolev-identity-inclusion-bounded}
Corollary~\ref{cor:sobolev-extension-continuity} 表明, 若 $k>m/2$, $\Omega\subset\mathbb{R}^m$ 具有 $k$ 阶延拓性质, 则从 $H^k(\Omega)$ 到 $C_b(\overline\Omega)$ 的恒等映射有界.
\end{Example}

\Needspace{4\baselineskip}
\begin{Example}
\label{ex:riesz-map-conjugate-linear-operator}
Riesz 表示定理断言, 每个 Hilbert 空间 $H$ 都与其对偶空间 $H^*$ 等距. 在 Remark~\ref{rem:riesz-mappings} 中构造了映射 $A_H:H^*\to H$, 它对每个线性泛函 $l\in H^*$ 给出唯一元素 $x=A_H(l)\in H$, 使\[l(y)=(x,y)\quad\text{for all }y\in H.\]注意, 算子 $A_H$ 实际上是\emph{共轭线性的}; 即 $A_H(\alpha l)=\bar\alpha A_H(l)$. 该算子有界; Riesz 表示定理保证其范数为 $1$.
\end{Example}

\Needspace{4\baselineskip}
\begin{Example}
\label{ex:right-left-shifts-domain-dependent-inverses}
对任意 $x=\{x_1,x_2,\ldots,x_n,\ldots\}\in\ell^2$, 用下式定义\textbf{右移}算子 $S_r:\ell^2\to\ell^2$:\[S_r(x)=\{0,x_1,x_2,\ldots,x_n,\ldots\}.\]值域为 $\ell^2$ 中第一分量为 $0$ 的所有元素的集合, 即 $\{1,0,0,\ldots\}^{\perp}$. 零空间为单点集 $\{0\}$. 由于 $\|S_r(x)\|=\|x\|$, 有 $\|S_r\|=1$.

用下式定义\textbf{左移}算子 $S_l:\ell^2\to\ell^2$:\[S_l(x)=\{x_2,\ldots,x_n,\ldots\}.\]$S_l$ 的值域为 $\ell^2$, 零空间由元素 $\{1,0,0,\ldots\}$ 张成. 由于 $\|S_l(x)\|\leq\|x\|$, 并且对任意 $x\in\{1,0,0,\ldots\}^{\perp}$ 都有 $\|S_l(x)\|=\|x\|$, 所以 $\|S_l\|=1$.

显然, $S_r$ 可逆, 且 $S_r^{-1}=S_l$. 若取 $\mathcal{D}(S_l)=\ell^2$, 则 $S_l$ 不可逆. (其零空间非平凡.) 然而, 若取 $\mathcal{D}(S_l):=\{1,0,0,\ldots\}^{\perp}$, 则 $S_l$ 可逆, 且 $S_l^{-1}=S_r$.
\end{Example}

\Needspace{4\baselineskip}
\begin{Example}
\label{ex:orthogonal-projection-operator}
设 $H$ 为 Hilbert 空间, $M\subset H$ 为闭子空间. 定义\textbf{正交投影}算子 $P_M:H\to M$, 使对任意 $x\in H$, $P_Mx$ 为 $M$ 中满足 $(x-P_Mx)\in M^{\perp}$ 的唯一元素.
\end{Example}
